{\tamilfont
சாலைகளில் மின்சார வாகனங்களின் (EVs) மொத்த எண்ணிக்கை அதிகரித்து வருவதால், அவற்றின் பேட்டரிகளின் பாதுகாப்பு மற்றும் நம்பகத்தன்மைக்கான தேவை மிகவும் முக்கியத்துவம் பெறுகிறது. டென்ட்ரைட் உருவாக்கம் என்பது ஒரு மின்சார வாகனத்தின் பேட்டரி எதிர்கொள்ளும் மிக முக்கியமான பிரச்சனைகளில் ஒன்றாகும். இது பாதுகாப்பு அபாயங்களை ஏற்படுத்துவதோடு மட்டுமல்லாமல், பேட்டரியையே அழித்துவிடவும் கூடும். இந்தத் திட்டம், பாரம்பரிய மற்றும் குவாண்டம் கணினி முறைகளை ஒன்றிணைக்கும் ஒரு கலப்பின இயந்திர கற்றல் மாதிரியைப் பயன்படுத்தி, பேட்டரி டென்ட்ரைட் அபாயத்தை முன்கூட்டியே கண்டறிவதற்கான ஒரு தனித்துவமான செயல்முறையை வழங்குகிறது. குவாண்டம் சப்போர்ட் வெக்டர் மெஷின்கள் (QSVM), மின்னழுத்தம், மின்னோட்டம் மற்றும் வெப்பநிலை ஆகியவற்றைக் கொண்ட பேட்டரியின் நேரத் தொடர் தரவுகளின் நிலையை வகைப்படுத்துகின்றன. பேட்டரியில் பொருத்தப்பட்டுள்ள சென்சார்கள் ஒரு குறிப்பிட்ட காலத்திற்கு தொடர்ச்சியான தரவுகளைச் சேகரிக்கின்றன. பின்னர் இந்தத் தரவு, ஸ்லைடிங் விண்டோ நுட்பத்தைப் பயன்படுத்தி சிறிய பகுதிகளாகப் பிரிக்கப்படுகிறது. ஒவ்வொரு பகுதிக்கும், பேட்டரியின் செயல்பாட்டில் உள்ள சாத்தியமான அம்சங்களைக் கண்டறிய, சராசரி, மாறுபாடு மற்றும் மாற்றத்தின் விகிதம் உள்ளிட்ட பல்வேறு புள்ளிவிவரங்கள் கணக்கிடப்படுகின்றன. தேர்ந்தெடுக்கப்பட்ட அம்சங்கள் பின்னர் கலப்பின மாதிரிக்கு வழங்கப்படுகின்றன. அங்கு QSVM, அம்ச உருமாற்றச் செயல்முறையை மேலும் முன்னெடுத்து, டென்ட்ரைட் வளர்ச்சியின் அறிகுறிகளான மிகச் சிறிய அசாதாரணங்களைக் கண்டறிவதை எளிதாக்குகிறது. டென்ட்ரைட் அபாயம் குறித்த நிகழ்நேர, மிகவும் துல்லியமான கணிப்புகளை உருவாக்குவதே இந்த மாதிரியின் நோக்கமாகும். இது விரைவான நடவடிக்கைகளை எடுக்கவும், அதன் மூலம் ஒட்டுமொத்த பேட்டரியை சிறப்பாக நிர்வகிக்கவும் வழிவகுக்கும். குவாண்டம் கணினி மற்றும் பாரம்பரிய இயந்திர கற்றல் நுட்பங்களின் இந்த இணைப்பே, மின்சார வாகனங்களின் பேட்டரி ஆரோக்கிய கண்காணிப்புத் துறையில் ஒரு முக்கிய மைல்கல்லாக அமைகிறது.
}